\pdfvariable suppressoptionalinfo 611
\providecommand{\CRevisionRound}{2}
\documentclass[11pt]{article}
\usepackage[a4paper,margin=25mm]{geometry}
\usepackage{amsmath,amssymb,amsthm,booktabs,microtype}
\usepackage[T1]{fontenc}
\usepackage{lmodern}
\usepackage[hidelinks]{hyperref}
\newtheorem{theorem}{Theorem}
\newtheorem{lemma}{Lemma}
\newcommand{\LV}{Lotka--Volterra}
\newcommand{\R}{\mathbb R}
\title{A Global Hamiltonian Period Atlas for Positive Lotka--Volterra Flow}
\author{HCS Research Program}
\date{28 August 2026\quad (revision \CRevisionRound)}
\begin{document}
\maketitle
\begin{abstract}
We give a global, source-native description of the positive two-species
Lotka--Volterra flow at its physical time clock.  A logarithmic normalization
turns the system into a strictly convex proper Hamiltonian flow.  Every
positive energy is therefore one compact periodic oval.  The two real
Lambert-W branches of the exponential potential give exact one-dimensional
quadratures for the enclosed area and period, and coarea gives the action
identity $J'(h)=T(h)/(2\pi)$.  Linearization supplies the center-period limit,
while integrating the log equations gives exact prey and predator cycle
averages.  An independent certificate checks 24 levels in six rational
parameter sets and re-integrates the period with a SciPy event ODE.  We do not
assert period monotonicity or a high-energy asymptotic, and we make no target
arithmetic or operator claim.
\end{abstract}

\section{Model and global claim}
Consider
\[
 \dot x=x(a-by),\qquad \dot y=y(-c+dx),\qquad a,b,c,d>0,
\]
on $x,y>0$, with physical time $t$.  Put $x_*=c/d$, $y_*=a/b$ and
$u=\log(x/x_*),v=\log(y/y_*)$.  Then
\[
 \dot u=a(1-e^v),\qquad \dot v=c(e^u-1),\qquad
 H=c(e^u-u-1)+a(e^v-v-1).
\]

\begin{theorem}[period annulus and exact atlas]\label{thm:atlas}
$H$ is a strict convex proper first integral with unique critical point
$(0,0)$.  Every level $H=h>0$ is a smooth compact periodic oval.  If
$F(s)=e^s-s-1$, let
\[
 \ell(r)=-W_0(-e^{-1-r})-1-r,\qquad
 q(r)=-W_{-1}(-e^{-1-r})-1-r.
\]
With $u_-=\ell(h/c)$, $u_+=q(h/c)$, and
$r(u)=(h-cF(u))/a$, its area and period are
\begin{align*}
 A(h)&=\int_{u_-}^{u_+}\bigl[q(r(u))-\ell(r(u))\bigr]\,du,\\
 T(h)&=\frac1a\int_{u_-}^{u_+}\left[
 \frac1{1-e^{\ell(r(u))}}+\frac1{e^{q(r(u))}-1}\right]\,du.
\end{align*}
The endpoint singularities are integrable.  For $J(h)=A(h)/(2\pi)$,
$J'(h)=T(h)/(2\pi)$, and
$\lim_{h\downarrow0}T(h)=2\pi/\sqrt{ac}$.  Every periodic orbit obeys
$\langle x\rangle=c/d$ and $\langle y\rangle=a/b$.
\end{theorem}

\begin{proof}
The Hessian is $\operatorname{diag}(ce^u,ae^v)$ and $F$ tends to infinity at
both ends, giving strict convexity and properness.  The sublevel
$K_h=\{H\le h\}$ is a compact convex body with nonempty interior, so its
boundary $H=h$ is connected; regularity follows because the only critical
point is the origin.  The canonical equations are
$\dot u=-H_v,\dot v=H_u$, and the vector field is nonzero on this boundary.
Thus the connected compact one-manifold is one periodic oval.  Solving
$F(s)=r$ gives the two displayed Lambert branches.  Splitting the oval into
lower and upper graphs gives the quadratures.  Along the level,
$|\dot{(u,v)}|=|\nabla H|$, hence $dt=ds/|\nabla H|$; the coarea formula
then gives $dA/dh=\int_{H=h}ds/|\nabla H|=T$.  The center matrix is
$\left(\begin{smallmatrix}0&-a\\c&0\end{smallmatrix}\right)$.
Finally, integrating $\dot u$ and $\dot v$ over a period gives
$\langle e^u\rangle=\langle e^v\rangle=1$.
\end{proof}

\ifnum\CRevisionRound>0
\section{Branches, action, and boundaries}
The branch labels matter: $W_0$ gives the negative graph and $W_{-1}$ the
positive graph.  A trigonometric substitution in the certificate cancels the
square-root endpoint singularities without changing the physical clock.  The
geometric action is positive by definition; no orientation convention is
needed.  The axes are invariant one-dimensional flows, not members of the
positive period annulus.  If a rate is zero, coercivity or the interior
Hamiltonian statement can fail and the resulting boundary equation is treated
separately.  We intentionally leave period monotonicity and large-energy
asymptotics outside the theorem.
\fi

\ifnum\CRevisionRound>1
\section{Independent audit and Route-A boundary}
The released ledger contains six positive rational $(a,b,c,d)$ cases and four
energies per case.  It serializes 24 areas, periods, actions, branch endpoints,
and residuals.  The producer uses high-precision Lambert-W quadrature; the
checker independently repeats those identities and then integrates the log
flow with a DOP853 event solver from the right turning point.  The checker
passes 732 assertions, including all 24 heterogeneous center-period checks;
SymPy passes 12 identities, byte replay is exact, and twelve repaired/stale
hash attacks are rejected.

This is a continuum of real-energy cycles, not a discrete arithmetic primitive
ledger.  Accordingly the Route-A tuple is
\begin{center}\small\ttfamily
(A0\_FAIL, A1\_WEAK, A2\_FAIL, A3\_FAIL, A4\_FORMAL\_HINT),
\end{center}
with \texttt{overall=ROUTE\_A\_REJECTED} and
\texttt{route\_b\_invocation\_allowed=false}.  No target prime, zero,
local factor, root number, automorphy statement, target functional equation,
or Hilbert--Polya operator is claimed.
\fi

\section*{Source note}
Period-analysis context includes Waldvogel~\cite{waldvogel1983, waldvogel1986}
and Hsu~\cite{hsu1983}.  We do not use the historical monotonicity result as a
claim here and make no priority statement.
\begin{thebibliography}{9}
\bibitem{waldvogel1983} J. Waldvogel, ``The Period in the Volterra--Lotka
Predator-Prey Model,'' \emph{SIAM Journal on Numerical Analysis} 20(6)
(1983), 1264--1272, DOI:\href{https://doi.org/10.1137/0720098}{10.1137/0720098}.
\bibitem{waldvogel1986} J. Waldvogel, ``The period in the Volterra-Lotka system
is monotonic,'' \emph{Journal of Mathematical Analysis and Applications} 114(1)
(1986), 178--184, DOI:\href{https://doi.org/10.1016/0022-247X(86)90076-4}{10.1016/0022-247X(86)90076-4}.
\bibitem{hsu1983} S.-B. Hsu, ``A Remark on the Period of the Periodic Solution
in the Lotka-Volterra System,'' \emph{Journal of Mathematical Analysis and
Applications} 95(2) (1983), 428--436,
DOI:\href{https://doi.org/10.1016/0022-247X(83)90117-8}{10.1016/0022-247X(83)90117-8}.
\end{thebibliography}

\section*{Declarations}
\textbf{Scope.} The exact registered literal is
\texttt{NO\_BAD\_EULER\_OR\_ROOT\_NUMBER}.
\textbf{Data and code.} The theorem, ledger, independent checks, and build
instructions are released with HCS-C211.
\textbf{Competing interests.} None declared.
\textbf{AI-use disclosure.} Generative tools assisted drafting and code
generation.  The released theorem, numerical values, source metadata, and
scope boundaries were checked by the internal artifact chain; this is not
external peer review.
\end{document}
